\documentclass[a4paper]{article}
\usepackage[pdftex]{graphicx}
\usepackage[utf8]{inputenc}
\usepackage{enumerate}
\usepackage{siunitx}
\sisetup{locale=DE} 
\DeclareSIUnit \parsec {pc}
\usepackage{href-ul}
\hypersetup{
	colorlinks=true,
	linkcolor=blue,
	urlcolor=blue}
\usepackage{icomma}
\usepackage{amssymb}
\usepackage{geometry}
\geometry{a4paper, top=15mm, left=15mm, right=15mm, bottom=15mm,
	headsep=10mm, footskip=12mm}
\hyphenation{Bahn-ra-dius}

%Source: img_0123

% Start the document
\begin{document}
%Titel
{\bf Schulaufgabe aus der Physik }\\
\begin{enumerate}[1.]

%Aufgabe 1
{\bf \item Der Planet Neptun hat acht bekannte Monde, von denen nur die beiden äußeren, Triton und Nereide, von der Erde aus entdeckt wurden.} Der 1846 vom englischen Amateurastronomen William Lassell beobachtete größte Mond Triton braucht für eine Umkreisung $T_T=5,88$ Tage bei einem mittleren Bahnradius von $r_T= \SI{354000}{\kilo\meter}$. \\
Nereide, 1949 vom niederländischen Astronomen Gerard Kuiper entdeckt, besitzt eine Umlaufdauer von $T_N=360,1 $ Tagen

\begin{enumerate}[a)]
\item Berechnen Sie die mittlere Bahngeschwindigkeit des Mondes Triton in $\SI{}{\kilo\meter\over \second}$
\item Ermitteln Sie den mittleren Bahnradius von Nereide.
\end{enumerate}

%Aufgabe 2
{\bf \item Erläutern Sie unter Verwendung des 2.Kepler'schen Gesetzes und mit Verwendung der Fachbegriffe Perihel und Aphel,}  warum das Winterhalbjahr auf der Nordhalbkugel um einige Tage kürzer ist als das Sommerhalbjahr.

%Aufgabe 3
{\bf \item Nennen Sie alle Ihnen bekannte astronomische Beobachtungsergebnisse, die als Hinweis für die Richtigkeit der Urknall-Theorie gelten.}

%Aufgabe 4
{\bf \item Die im Jahre 1789 vom deutschen Astronom Wilhelm Herschel entdeckte Galaxie NGC 3738} entfernt sich mit einer Geschwindigkeit von $v_G=\SI{331}{\kilo\meter\over \second}$ von uns. Berechnen Sie die Entfernung dieser Galaxie in Mpc und in km.\\
$[\textnormal{Hubble-Konstante:}~H_0=\SI{74}{\kilo\meter\over \second\cdot\mega\parsec}; \quad 1~\textnormal{Parsec~pc}=\SI{3,1e13}{\kilo\meter}]$

%Aufgabe 5
{\bf \item Der Dreiecksnebel M 33 gehört zur lokalen Nebelgruppe der Milchstraße.} Er hat eine Entfernung von $3,0 \cdot 10^6 $ Lichtjahren zu ihr und bewegt sich mit der Geschwindigkeit von $\SI{190}{\kilo\meter\over \second}$ auf sie zu. 

\begin{enumerate}[a)]
\item In wievielen Jahren werden sich die Milchstraße und M 33 vermischen?
\item Geben Sie einen Grund an, warum sich M 33 nicht von uns fortbewegt.
\end{enumerate} 

%Aufgabe 6
{\bf \item In nebenstehender Graphik ist die schematische Seitenansicht unserer Galaxie, der Milchstraße, dargestellt.}\\

\begin{minipage}{0.4\textwidth}
Beschriften Sie sie mit folgenden Fachbegriffen:\\

 Diskusscheibe\\
 Halo \\
 Staubschicht \\
 Zentrale Aufwölbung \\
 Kern \\
 Position Sonne \\
 Kugelsternhaufen. 
\end{minipage}
\hfill
\begin{minipage}{0.55\textwidth}
\includegraphics[width=7 cm]{milchstrasse123.pdf}
\end{minipage}


%Aufgabe 7
{\bf \item Korrigieren Sie alle Fehler in den folgenden Aussagen:}

\begin{enumerate}[a)]
\item Planeten bewegen sich auf Kreisbahnen um die Sonne, die im Mittelpunkt steht.
\vspace{35 pt}\\
\item Die Sonne, das Zentrum unseres Planetensystems, bewegt sich nicht.
\vspace{35 pt}\\
\item In gleichen Zeitabständen legt ein Planet bei seinem Umlauf um die Sonne die gleichen Strecken zurück.
\vspace{35 pt}\\
\item Befindet die Erde sich im Aphel ihrer Bahn, dann hat sie den kleinsten Abstand zur Sonne. 
\vspace{35 pt}\\
\item Der amerikanische Astronom Edwin  Hubble entdeckte eine Blauverschiebung im Spektrum der Sterne und folgerte, dass alle Galaxien den gleichen Abstand zur Erde haben.
\vspace{35 pt}\\
\item Wenn ein Radfahrer mit $\SI{18}{\kilo\meter\over \hour} $ fährt, dann legt er in 10 Minuten $\SI{1,8}{\kilo\meter}$ zurück.
\vspace{35 pt}\\
\item Im Mittelpunkt der Milchstraße, die im Universum still steht, befindet sich unsere Sonne. 
\end{enumerate} 

\end{enumerate}

\fbox{
	\begin{minipage}{0.5\textwidth}
		Zur Lösung bitte \href{https://www.okuyakl.de/phys/p10astL123/ll123.pdf}{hier klicken} oder den QR-Code scannen.\\
	Weitere Arbeitsblätter gibt es unter 
	
	\href{https://www.okuyakl.de}{www.okuyakl.de}
	\end{minipage}
	\hfill
	\begin{minipage}{0.4\textwidth}
		\includegraphics[width=1.5 cm]{../../viecher/zwe03}
		\includegraphics[width=3 cm]{qr123}
		\includegraphics[width=2 cm]{../../viecher/afanticon1}
		
	\end{minipage}}

\end{document}%L LÖSUNG ---------------------------------------------------
